\documentclass{report}
\usepackage{amsmath,amssymb}
\setlength{\parindent}{0mm}
\setlength{\parskip}{1em}
\begin{document}
\begin{center}
\rule{6in}{1pt} \
{\large Jon Calhoun \\
{\bf Understanding the Propagation of Silent Data Corruption in Algebraic Multigrid}}

201 North Goodwin Ave \\ Urbana \\ IL 61802
\\
{\tt jccalho2@illinois.edu}\\
Luke Olson\end{center}

Sparse linear solvers from a fundamental kernel in high performance computing (HPC).
Exascale systems are expected to be more complex than systems of today
being composed of thousands of
heterogeneous processing elements that operate at near-threshold-voltage to meet
power constraints. The combination of near near-threshold-voltage and number of
processing elements required to reach exascale increases the rate of
silent data corruption (SDC).
With the rate of SDC expected
to be higher, understanding how error propagates in HPC applications becomes vital
to devise efficient detection and recovery schemes. In this talk, we investigate
how SDC occurring in fixed-point and floating-point instructions
propagates in the linear solver algebraic multigrid (AMG). We discover
that
SDC occurring on the coarsest levels have the most impact on convergence
requiring extra iterations in a higher percentage than on the finest
levels.


\end{document}
